\documentclass{article}

\usepackage{amsmath,amssymb}
\usepackage{xurl}
\usepackage[hidelinks]{hyperref}
\setlength{\emergencystretch}{2em}

% Shared commands for notes in docs/paper-gaps/.

\DeclareUnicodeCharacter{2080}{\ifmmode{}_{0}\else\textsubscript{0}\fi}
\DeclareUnicodeCharacter{2081}{\ifmmode{}_{1}\else\textsubscript{1}\fi}
\DeclareUnicodeCharacter{2082}{\ifmmode{}_{2}\else\textsubscript{2}\fi}
\DeclareUnicodeCharacter{2099}{\ifmmode{}_{n}\else\textsubscript{n}\fi}

\newcommand{\Lean}{\textsc{Lean}}
\ExplSyntaxOn
\NewDocumentCommand{\leanid}{m}{
  \ifmmode
    \textsf{\detokenize{#1}}
  \else
    \group_begin:
    \urlstyle{sf}
    \tl_set:Nn \l_tmpa_tl {#1}
    \tl_replace_all:Nnn \l_tmpa_tl {\_} {_}
    \exp_args:NV \path \l_tmpa_tl
    \group_end:
  \fi
}
\ExplSyntaxOff
\providecommand{\tr}{\operatorname{tr}}
\newcommand{\tnleanIssueBase}{https://github.com/LionSR/TNLean/issues/}
\newcommand{\tnleanPrBase}{https://github.com/LionSR/TNLean/pull/}
\newcommand{\tnleanDocBase}{https://sirui-lu.com/TNLean/docs/}
\newcommand{\ghissue}[1]{\href{\tnleanIssueBase#1}{GitHub issue \##1}}
\newcommand{\ghpr}[1]{\href{\tnleanPrBase#1}{PR \##1}}
\newcommand{\doclink}[1]{\href{\tnleanDocBase#1.html}{\path{#1}}}

% Dirac notation and the identity operator, as in blueprint/src/macros/common.tex.
% \providecommand keeps note-local definitions authoritative.
\usepackage{dsfont}
\providecommand{\ket}[1]{|#1\rangle}
\providecommand{\bra}[1]{\langle #1|}
\providecommand{\braket}[2]{\langle #1 | #2 \rangle}
\providecommand{\ketbra}[2]{|#1\rangle\!\langle #2|}
\providecommand{\Id}{\mathds{1}}
\providecommand{\one}{\mathds{1}}
\providecommand{\C}{\mathbb{C}}
\providecommand{\MN}[1]{M_{#1}(\C)}
\providecommand{\MD}{M_D(\C)}

% Verdict marker: \gapnote{<kind>}{<status>}. Kind is the policy
% classification (clarification, local-correction, scope-restriction,
% unfaithful, false-source, open-gap); status is open, wip (elimination
% underway), resolved, or historical. Machine-read by the site index;
% prints a verdict line.
\newcommand{\gapnote}[2]{\par\noindent\textbf{Verdict:} #1 (#2).\par}


% Fallbacks keep this note compilable when a different command.tex is found first.
\providecommand{\leanid}[1]{\textsf{\detokenize{#1}}}
\providecommand{\gapnote}[2]{\par\noindent\textbf{Verdict:} #1 (#2).\par}

\title{Party-Network Bounds in the Dyadic Routing Lemma}
\author{}
\date{2026-10-08}

\begin{document}
\maketitle
\gapnote{scope-restriction}{open}

\section{The source lemma}

Lemma~8.2 of the polynomial PEPS approximation manuscript of September~24,
2026~\cite{OpenAI2026PolynomialPEPS} (\path{07-assembly.tex},
lines~122--127) states, with no hypotheses, that the links of the party
network can be routed on the padded grid $\{0,\ldots,N-1\}^2$ with constant
edge congestion, and that after reflection into the genuine $L\times L$ grid
the congestion is still constant and the physical output coordinates are
unchanged. The lemma concerns one specific network: the party network built
by the distribution protocol of Proposition~7.1 (\path{06-geometry.tex},
lines~12--53), with its parties attached to the anchors of the dyadic blocks
(\path{07-assembly.tex}, lines~103--114). The proof (lines~128--162) uses
four properties of that network:
\begin{enumerate}
\item every anchor lies on the padded grid (lines~113--114);
\item at most a bounded number $B$ of parties is attached to any dyadic block
  (lines~113--114, from item~1 of Proposition~7.1);
\item every party has at most a bounded number $\Delta$ of incident links
  (lines~140--141, from item~2 of Proposition~7.1);
\item the anchors of the two endpoints of every link are at distance at most
  $c$ times either endpoint scale, for a fixed $c$ (lines~130--132, from
  item~3 of Proposition~7.1).
\end{enumerate}
Proposition~7.1 also gives that the two endpoint scales of every link are
equal or adjacent; the congestion bound does not use this.

\section{The formal statement}

Proposition~7.1 is not yet formalized, so no formal party network of the
protocol exists to which the lemma could refer. The formal statements are
therefore made for an arbitrary party network with parties attached to
dyadic blocks, and the four properties above are hypotheses: the anchors lie
on the padded grid, at most $B$ parties per block, at most $\Delta$ links per
party, and link separation at most $c$ times either endpoint scale. Under
these hypotheses the padded congestion is at most $2(2c+1)B\Delta$, the
folded congestion is at most $8(2c+1)B\Delta$, and the routed network is
exactly a PEPS on the open $L\times L$ square lattice with bond dimension at
most $\max(D,1)^{8(2c+1)B\Delta}$ when every link has dimension at most
$D$.\footnote{The padded bounds are
\leanid{TNLean.PEPS.Approximation.DyadicPlacement.load_horizontal_le} and
\leanid{TNLean.PEPS.Approximation.DyadicPlacement.load_vertical_le}; the
folded bound is
\leanid{TNLean.PEPS.Approximation.foldedRouting_card_traversal_le}; the PEPS
assembly is
\leanid{TNLean.PEPS.Approximation.exists_squarePEPS_of_dyadicRouting}.}

The constants depend only on $B$, $\Delta$ and $c$, as in the source. The
formal statements are a conditional form of Lemma~8.2, valid for every party
network with the four properties; the lemma for the network of the source is
the special case in which the four properties hold.

\section{Elimination}

Formalize Proposition~7.1 and the attachment of its parties to the anchors
of the dyadic blocks of the padded hierarchy. The two sources of the four
properties are then discharged separately. Property~1, that every anchor
lies on the padded grid, follows from the attachment itself
(\path{07-assembly.tex}, lines~113--114) and is not a clause of
Proposition~7.1. Properties~2--4, the block, degree and separation bounds,
follow from items~1--3 of Proposition~7.1, with constants uniform in the
Hamiltonian and in $L$. Applying the conditional statements to that network
then gives Lemma~8.2 as stated.

\bibliographystyle{alpha}
\bibliography{references}

\end{document}
